%\ifodd\value{page}\hbox{}\newpage\fi
\chapter*{Śrī Lalitā Sahasranāma — Śloka 2}
\addcontentsline{toc}{chapter}{Śloka 2 - Nomi 4,5}

\begin{sloka}
{\sanskritfont
\begin{tabular}{c}
उद्यद्भानुसहस्राभा चतुर्बाहुसमन्विता ।\\
रागस्वरूपपाशाढ्या क्रोधाकाराङ्कुशोज्ज्वला ॥ २ ॥
\end{tabular}

\vspace{1.5em}

{\middlesize \iastfont
udyadbhānu­sahasrābhā caturbāhu­samanvitā |\\
rāga­svarūpa­pāśāḍhyā krodhākārāṅkuśojjvalā ||2||}

\vspace{1em}

{\middlesize
\anvaya{%
उद्यद्भानुसहस्राभा चतुर्बाहुसमन्विता
रागस्वरूपपाशाढ्या क्रोधाकाराङ्कुशोज्ज्वला ।
}
}

\middlesize \iastfont \itmean{%
«Splendente come mille soli sorgenti, dotata di quattro braccia;  
che impugna il laccio la cui essenza è l'attrazione  
e il pungolo che risplende come forma del controllo.»%
}

}
\end{sloka}

\wordmeaning{%
\wordline{उद्यद्-भानु-सहस्र-आभा}{udyad-bhānu-sahasra-ābhā}
  { splendente come mille soli sorgenti;}%
\wordline{चतुर्-बाहु-समन्विता}{catur-bāhu-samanvitā}
  { dotata di quattro braccia;}%
\wordline{राग-स्वरूप-पाश-आढ्या}{rāga-svarūpa-pāśāḍhyā}
  { che possiede il laccio la cui essenza è l’attrazione (\textit{rāga});}%
\wordline{क्रोध-आकार-अङ्कुश-उज्ज्वला}{krodha-ākāra-aṅkuśa-ujjvalā}
  { che risplende con il pungolo nella forma del controllo (\textit{krodha}).}%
}

Lo \textit{Śloka} 2 sviluppa l’aspetto dinamico della sovranità enunciata nello \textit{śloka} precedente, mostrando \textit{Lalitā} come potenza luminosa in atto. L’espressione \textit{udyadbhānu-sahasrābhā} non descrive una luce statica, ma una luce sorgente, indicativa di una coscienza che emerge e opera. La qualifica \textit{caturbāhu-samanvitā} segnala una capacità di azione ordinata su più livelli dell’esperienza. 

Il \textit{pāśa} e l’\textit{aṅkuśa} non sono presentati come armi materiali, ma come strumenti funzionali dell’azione divina: il \textit{pāśa}, detto \textit{rāga-svarūpa}, è la forza di attrazione mediante cui la coscienza lega le forme a sé e le mantiene nella relazione, rendendo possibile la coesione della manifestazione; l’\textit{aṅkuśa}, qualificato come \textit{krodha-ākāra}, è la forza direttiva che arresta la deviazione e orienta il movimento, impedendo che l’attrazione si trasformi in dispersione. Desiderio e controllo non sono quindi passioni negate, ma energie assunte e governate dalla coscienza divina, attraverso le quali l’azione cosmica di \textit{Lalitā} si compie in modo ordinato, consapevole e luminoso.

Lo \textit{Śloka} 2 contiene due \textit{nāma} propriamente detti, che descrivono la modalità dell’azione divina. \textbf{\textit{Udyadbhānu-sahasrābhā}} definisce \textit{Lalitā} come potenza luminosa sorgente, indicando una coscienza in atto e non statica; \textbf{\textit{caturbāhu-samanvitā}} esprime la capacità di operare simultaneamente su più livelli. Le espressioni \textit{rāga-svarūpa-pāśāḍhyā} e \textit{krodhākārāṅkuśojjvalā} non costituiscono nuovi \textit{nāma} autonomi, ma specificano gli strumenti e le energie attraverso cui questi due nomi operano, qualificando l’attrazione e il controllo come funzioni integrate della coscienza divina.
